%!TEX root = ../main.tex
% =============================================================
%  CAPÍTULO 4 — MARCO TECNOLÓGICO       (Entrega 4 · 10 % · 25/8)
% =============================================================

\chapter{Marco tecnológico}\label{cap:marco-tecnologico}

\begin{guia}[Guía de la cátedra: qué se evalúa en este capítulo]
Documenta el \emph{cómo}. No es un manual de uso: es la justificación
de cada elección tecnológica. Regla: nada ``porque sí''.
\begin{enumerate}
    \item \textbf{Arquitectura}: diagrama con componentes y conexiones,
          responsabilidad y tecnología de cada uno, flujo de datos y
          justificación del estilo arquitectónico con alternativas
          evaluadas \citep{sommerville2016}.
    \item \textbf{Prototipo o \gls{g:mvp}}: cuál es y por qué; alcance
          implementado vs.\ simulado; capturas o enlace; cómo se probó.
    \item \textbf{Stack}: por tecnología, nombre, versión, capa, por qué
          (madurez, comunidad, costo, conocimiento del equipo) y qué
          alternativas se descartaron.
    \item \textbf{Infraestructura y base de datos}: dónde se aloja cada
          componente, escalabilidad, costos (enlazados al
          capítulo~\ref{cap:economica}), motor y modelo de datos.
    \item \textbf{Diseño y \gls{ux}}: quién lo usa, mockups, decisiones de
          diseño y \textbf{cómo se validó con usuarios} (la cátedra lo marca
          como importante; cuestionario en el apéndice~\ref{apx:encuesta}).
    \item \textbf{Seguridad, calidad y despliegue}: autenticación, control
          de acceso, cifrado, plan de pruebas, repositorio, \gls{cicd} y
          monitoreo.
\end{enumerate}
\end{guia}

\section{Introducción}

Este capítulo explica cómo se lleva la propuesta planteada en los
capítulos anteriores a un sistema real y qué tecnologías se utilizan
para desarrollarlo. El proyecto se desarrolló de manera progresiva:
la primera etapa fue un sistema de gestión aplicado a Rodríguez
Catering, que permite administrar clientes, eventos, presupuestos,
contratos, productos y proveedores; sobre esa misma base se incorporaron
después el esquema multi-tenant, el marketplace, el sistema de
recomendación mediante \gls{ia} y el Gemelo Digital del evento.

\section{Alcance del prototipo}\label{sec:alcance-prototipo}

El sistema desarrollado permite gestionar clientes, eventos,
presupuestos, contratos, productos y proveedores de forma relacionada: un
cliente puede solicitar un presupuesto y confirmar un evento, y ese
evento puede necesitar distintos productos y servicios. Sobre esa base se
implementaron las siguientes capas, en este orden:

\begin{enumerate}
    \item \textbf{Esquema multi-tenant}: las tablas \texttt{salones} y
          \texttt{proveedores\_servicios} representan a los distintos
          tenants que pueden ofrecer servicios en el marketplace (Rodríguez
          Catering quedó registrado como el primer tenant), y
          \texttt{reservas} es la tabla de unión entre un evento y el
          tenant contratado.
    \item \textbf{Autenticación real}: se reemplazó el token falso en
          base64 del prototipo original por Laravel Sanctum, con roles
          (cliente, salón, proveedor, admin) y activación del portal por
          invitación.
    \item \textbf{Marketplace público}: búsqueda y ficha de salones y
          proveedores con disponibilidad, y solicitud de contratación que
          genera el evento y la reserva correspondiente, reemplazando el
          flujo anterior por WhatsApp.
    \item \textbf{Motor de recomendación}: aislado en
          \texttt{App\textbackslash Services\textbackslash Recommendation},
          con scoring por reglas ponderadas (sección~\ref{sec:motor-recomendacion}).
    \item \textbf{Gemelo Digital}: simulador de escenarios por evento
          (sección~\ref{sec:gemelo-digital}).
\end{enumerate}

También se desarrolló un portal destinado a los clientes, que les permite
consultar el estado de un presupuesto o pedido sin comunicarse
constantemente con la empresa.

Quedan pendientes los pagos en sandbox (todavía no se eligió pasarela
entre Mercado Pago y Stripe test mode) y el despliegue público (todavía
no hay cuentas de hosting creadas); ambos se dejaron deliberadamente para
el final, ya que no condicionan la validación funcional del resto de la
plataforma.

\section{Arquitectura del sistema}

Se eligió una arquitectura cliente-servidor, separando el frontend
(pantallas con las que interactúan los usuarios) del backend (encargado
de procesar y almacenar la información). Existe un frontend para el
personal interno y un portal separado para los clientes externos, ambos
comunicándose con el mismo backend a través de una \gls{api} REST.

El funcionamiento general puede resumirse como: Usuario $\rightarrow$
Frontend $\rightarrow$ API REST $\rightarrow$ Lógica del sistema
$\rightarrow$ Base de datos. Esta estructura facilitó la incorporación
del motor de recomendación y del Gemelo Digital como nuevos componentes
que reciben información del sistema y devuelven sus resultados, sin
modificar la arquitectura existente: ambos quedaron implementados como
servicios de Laravel (\texttt{App\textbackslash Services\textbackslash
Recommendation} y \texttt{App\textbackslash Services\textbackslash
DigitalTwin}) inyectados en los controladores correspondientes, no como
sistemas aparte.

\begin{figure}[H]
    \centering
    \fbox{\parbox{0.85\textwidth}{\centering\vspace{2em}
    \completar{Diagrama de arquitectura cliente-servidor de la plataforma
    (frontend interno + portal de clientes $\to$ API REST $\to$ Laravel
    $\to$ MySQL/SQLite, con los servicios de Recomendación y Gemelo
    Digital colgando del backend).}
    \vspace{2em}}}
    \caption{Arquitectura cliente-servidor de la plataforma}
    \label{fig:arquitectura}
\end{figure}

\subsection{Justificación del tipo de arquitectura}

Se optó por una arquitectura cliente-servidor de tres capas en lugar de
una arquitectura de microservicios. En la etapa actual del proyecto, la
cantidad de funcionalidades y de usuarios concurrentes no justifica la
complejidad operativa adicional que implican los microservicios
(orquestación, comunicación entre servicios, despliegues independientes).
La arquitectura elegida permite que un único backend sirva tanto al panel
interno como al portal de clientes, evitando duplicar lógica de negocio,
y deja abierta la posibilidad de extraer componentes específicos ---como
el motor de recomendación--- como servicios independientes en etapas
posteriores, cuando el volumen de uso lo justifique. El motor de
recomendación ya se diseñó detrás de una interfaz
(\texttt{RecommendationEngine}) precisamente para poder reemplazar la
implementación por reglas por un modelo entrenado sin tocar los
controladores que lo consumen.

\section{Prototipo funcional}

Lo desarrollado no consiste en pantallas de ejemplo, sino en un sistema
funcional conectado a una base de datos: permite registrar un cliente y
utilizar esa información para generar presupuestos, eventos y contratos,
buscar y solicitar salones o proveedores en el marketplace, recibir
recomendaciones y simular escenarios de un evento con el Gemelo Digital.
Las pruebas realizadas hasta el momento se orientan principalmente a
verificar que las funcionalidades desarrolladas respondan correctamente;
a medida que avance el proyecto se realizarán pruebas más completas con
usuarios reales.

\section{Tecnologías utilizadas}

Para desarrollar la parte visual se utiliza React junto con TypeScript,
complementado con Material UI para los componentes visuales. Para el
backend se utiliza Laravel con PHP, junto con Eloquent para trabajar con
la base de datos. La información se almacena en una base de datos
relacional: MySQL para el entorno de producción y SQLite durante el
desarrollo.

\subsection{Alternativas evaluadas y descartadas}

Tal como pide la cátedra, se detalla a continuación, para cada
tecnología, la versión utilizada, las alternativas consideradas y el
motivo por el cual fueron descartadas.

\begin{table}[H]
    \centering
    \caption{Stack tecnológico y alternativas evaluadas}
    \label{tab:stack}
    \small
    \begin{tabularx}{\textwidth}{l l L L L}
        \toprule
        \tb{Tecnología} & \tb{Versión} & \tb{Uso} & \tb{Alternativas evaluadas} & \tb{Por qué se descartaron} \\
        \midrule
        React & 18.x & Frontend (panel interno y portal de clientes) & Vue.js, Angular & Vue se descartó por tener menor disponibilidad de componentes que React; Angular se descartó por su curva de aprendizaje más pronunciada y por ser más pesado de lo necesario para el alcance actual. \\
        TypeScript & 5.x & Tipado estático sobre JavaScript & JavaScript puro & Se descartó JavaScript puro porque, con muchas entidades relacionadas (clientes, eventos, presupuestos, productos), el tipado estático reduce errores en desarrollo y facilita el mantenimiento. \\
        Material UI & 5.x & Componentes visuales & Tailwind CSS, Chakra UI & Tailwind hubiera exigido construir los componentes desde más abajo; Chakra UI tenía comunidad y documentación más acotadas al iniciar el proyecto. \\
        Laravel (PHP) & 10.x & Backend / API REST & Node.js con Express, Django (Python) & Node.js se descartó porque Laravel integra autenticación, migraciones y ORM que en Node se debían incorporar por separado; Django se descartó por menor familiaridad previa con Python. \\
        Eloquent (ORM) & incluido en Laravel 10.x & Mapeo objeto-relacional & Query Builder sin ORM, Doctrine & Eloquent simplifica la definición de relaciones entre entidades; Doctrine no es el ORM nativo de Laravel y añadía complejidad de integración. \\
        MySQL / SQLite & MySQL 8.x (prod.) / SQLite (dev.) & Persistencia de datos & PostgreSQL, MongoDB (NoSQL) & PostgreSQL no aportaba ventajas relevantes frente a MySQL para este volumen de datos; MongoDB se descartó porque la información es altamente relacional. \\
        \bottomrule
    \end{tabularx}
\end{table}

\section{Infraestructura y alojamiento}

Durante la etapa de desarrollo, el sistema se ejecuta de manera local.
Cuando se publique, se necesitará un servidor que permita ejecutar el
backend en PHP/Laravel, conectarse a una base de datos MySQL y almacenar
los archivos que genere o reciba la aplicación. El frontend puede
alojarse en el mismo servidor o en un servicio diferente, dado que se
encuentra separado del backend. Si la plataforma crece, se podrá evaluar
una infraestructura en la nube que permita escalar según la cantidad de
usuarios. La elección concreta del proveedor de hosting es uno de los dos
puntos que se dejaron para el final del proyecto.

\section{Base de datos y modelo de información}

Entre los datos principales que se administran se encuentran clientes,
eventos, presupuestos, contratos, productos, proveedores, usuarios,
salones, proveedores de servicio y reservas, vinculados entre sí mediante
una base de datos relacional. Los cambios sobre la estructura de la base
de datos se administran mediante las migraciones que proporciona
Laravel, lo que permite registrar las modificaciones y recrear la misma
estructura en otro entorno cuando sea necesario.

\subsection{Estrategia de backups}

\begin{itemize}
    \item \textbf{Backups automáticos diarios}: copia completa de la
          base de datos una vez por día, en horario de bajo uso, mediante
          las herramientas del hosting o un paquete de Laravel como
          \texttt{spatie/laravel-backup}.
    \item \textbf{Retención}: conservar los últimos 7 backups diarios y
          un backup semanal durante las últimas 4 semanas.
    \item \textbf{Backups antes de migraciones}: backup manual adicional
          antes de aplicar cualquier migración que modifique tablas
          existentes en producción.
    \item \textbf{Almacenamiento externo}: guardar las copias en un
          almacenamiento distinto al del servidor de base de datos.
    \item \textbf{Prueba de restauración}: verificar periódicamente que
          un backup pueda restaurarse correctamente en un entorno de
          prueba.
\end{itemize}

\section{Diseño y experiencia de usuario}

El sistema interno concentra mayor cantidad de herramientas de gestión,
mientras que el portal del cliente se limita a las opciones que necesita,
evitando exponer funciones internas. Se busca mantener un mismo criterio
visual en todas las pantallas, apoyado en Material UI, y adaptar la
interfaz a computadoras, tablets y dispositivos móviles, dado que tanto
clientes como proveedores pueden necesitar consultar información desde
el celular.

\section{Seguridad, calidad y pruebas}

El sistema identifica quién está ingresando mediante Laravel Sanctum y,
según su rol (cliente, salón, proveedor, admin), determina qué
información puede consultar o modificar. Las contraseñas se almacenan de
forma segura, y la comunicación entre navegador y servidor deberá
realizarse mediante HTTPS una vez publicado el sistema. En cuanto a
pruebas, el objetivo actual es comprobar que las funcionalidades
respondan correctamente, con pruebas funcionales sobre el modelo de
autenticación y los flujos del portal de clientes; a medida que avance el
proyecto se automatizarán también las verificaciones sobre cálculos de
presupuestos, el motor de recomendación y el simulador del Gemelo
Digital.

\section{El marketplace, la recomendación y el Gemelo Digital}

Sobre la base del sistema de gestión y el esquema multi-tenant se
construyeron las tres piezas que diferencian a esta propuesta de un CRUD
convencional.

\subsection{Marketplace público}

El marketplace expone la búsqueda y la ficha de salones y proveedores
(con su disponibilidad) sin necesidad de iniciar sesión, y una solicitud
de contratación que, una vez autenticado el cliente, genera el evento y
la reserva correspondiente del lado del tenant elegido. Esto reemplaza el
flujo anterior, que dependía de WhatsApp para iniciar cualquier consulta.

\subsection{Motor de recomendación}\label{sec:motor-recomendacion}

El motor de recomendación está aislado detrás de una interfaz
(\texttt{RecommendationEngine}) con dos métodos, \texttt{recomendarSalones}
y \texttt{recomendarProveedores}, de forma que la implementación actual
pueda reemplazarse por un modelo entrenado más adelante sin modificar los
controladores que la consumen --- todavía no hay suficiente historial de
uso real para entrenar un modelo, así que la primera versión
(\texttt{RuleBasedRecommendationEngine}) puntúa cada salón o proveedor
con una combinación ponderada y configurable de:

\begin{itemize}
    \item \textbf{Presupuesto}: puntaje máximo si el presupuesto del
          evento cubre el rango de precios del tenant; se degrada de
          forma proporcional cuanto más caro sea el mínimo respecto al
          presupuesto.
    \item \textbf{Capacidad} (salones) o \textbf{ubicación}
          (proveedores): puntaje máximo si la cantidad de personas entra
          en el rango de capacidad del salón, o si la zona de cobertura
          del proveedor coincide con la ubicación del evento.
    \item \textbf{Reseñas}: puntaje proporcional al promedio de reseñas
          del tenant.
\end{itemize}

Antes de puntuar, se descartan directamente los salones o proveedores que
ya tienen una reserva confirmada para la misma fecha del evento. El
resultado incluye, además del puntaje, un motivo breve (por ejemplo,
``destaca por dentro de tu presupuesto y buenas reseñas de otros
clientes'') armado a partir de los componentes que superaron un umbral,
para que la recomendación sea explicable y no una caja negra.

\subsection{Gemelo Digital}\label{sec:gemelo-digital}

El \texttt{SimulationService} recibe un evento y un conjunto de cambios
hipotéticos (cantidad de invitados, fecha o presupuesto) y calcula el
impacto de esos cambios \emph{sin} modificar el evento real:

\begin{itemize}
    \item Reevalúa si la nueva cantidad de invitados sigue entrando en la
          capacidad del salón ya reservado (\emph{ok}, \emph{excede} o
          \emph{por debajo}) y recalcula la cantidad de mesas necesarias.
    \item Reestima el costo del catering: si ya hay una reserva de
          catering con un monto real cargado, escala ese monto
          proporcionalmente a la nueva cantidad de invitados; si no, usa
          un costo estimado por persona configurable.
    \item Suma el costo de las demás reservas confirmadas para obtener un
          presupuesto total estimado y marca si ese total supera el
          presupuesto declarado para el evento.
    \item Si la fecha cambia, revisa si alguno de los tenants ya
          confirmados para este evento tiene otra reserva confirmada
          justo en la nueva fecha, y devuelve el conflicto con su nombre
          comercial.
\end{itemize}

A diferencia del motor de recomendación, la cátedra no pide que el
Gemelo Digital sea reemplazable por un modelo de \gls{ml}, por lo que se
implementó como una única clase de servicio en lugar de una interfaz.
Cada simulación queda registrada, y el cliente puede aplicar una
simulación para que sus cambios pasen a ser el estado real del evento.

\section{Factibilidad técnica}

El sistema actual ya permite trabajar con clientes, eventos,
presupuestos, contratos, productos, proveedores, salones, proveedores de
servicio y reservas de forma relacionada, y ya tiene en funcionamiento el
marketplace, el motor de recomendación y el Gemelo Digital descriptos
más arriba. Esto confirma que la arquitectura elegida permitió incorporar
estos tres componentes sin reconstruir el sistema existente, tal como se
había previsto en la sección~\ref{sec:alcance-prototipo}.

\section{Limitaciones actuales y próximos pasos}

Quedan dos componentes pendientes, deliberadamente postergados al final
del proyecto porque no condicionan la validación funcional del resto de
la plataforma:

\begin{itemize}
    \item \textbf{Pagos en sandbox}: todavía no se eligió la pasarela
          (Mercado Pago o Stripe test mode).
    \item \textbf{Despliegue público}: todavía no existen las cuentas de
          hosting; el sistema corre únicamente en el entorno local de
          desarrollo.
\end{itemize}

El motor de recomendación, al estar basado en reglas y no en un modelo
entrenado, tiene una limitación conocida: no aprende del historial de
contrataciones ni de las valoraciones más allá del promedio de reseñas.
Reemplazarlo por un modelo entrenado, como trabajo futuro, requiere antes
acumular uso real de la plataforma, ya que hoy no hay suficiente
historial para entrenarlo.
